% =========================================================================
% Chapter 6 - Results
% =========================================================================
\chapter{Results}
\label{ch:results}

\section{The data set}

Three platinum structures, each a $3\times3\times3$ supercell of the face-centred
cubic primitive cell ($27$ atoms, $Z^{\mathrm{val}}=11$, so $297$ valence
electrons), computed with VASP~6.2.0 at
$E_{\mathrm{cut}}=450$\,eV, \opt{PREC=Accurate}, and a $5\times5\times5$
$\Gamma$-centred $k$-mesh.

\begin{center}
\begin{tabular}{@{}llrl@{}}
\toprule
Structure & Geometry & $\Omega$ [\AA$^3$] & Native grid \\
\midrule
struct\_000 & pristine fcc & $484.53$ & $128\times128\times128$ \\
struct\_001 & rattled $+$ strained & $477.65$ & $120\times128\times128$ \\
struct\_002 & rattled $+$ strained & $476.92$ & $128\times120\times128$ \\
\bottomrule
\end{tabular}
\end{center}

The three native grids \emph{differ}, which is not an inconvenience but the
central test: it is exactly the situation Section~\ref{sec:discinv} was designed
for, and it arises automatically because the cell shape enters VASP's grid
selection.

\begin{ptip}
Fields are reduced to a longest axis of $32$ points by the spectral truncation
of Section~\ref{sec:resample}, giving $32^3$, $30\times32\times32$ and
$32\times30\times32$. The shape difference \emph{survives} the reduction --- had
a fixed target shape been used instead, the ragged-grid capability would have
gone untested.
\end{ptip}

\section{Verification before learning}

Nothing was trained until the field machinery was checked against quantities
known independently.

\begin{center}
\begin{tabular}{@{}lll@{}}
\toprule
Check & Expected & Measured \\
\midrule
Poisson residual, $\nabla^2\vext$ vs $4\pi e^2 n_{\mathrm{ion}}$
  & $0$ & $4.5\times10^{-12}$ \\
Translation covariance & exact & $2.1\times10^{-14}$ \\
Linearity in $Z^{\mathrm{val}}$ & exact & $0.0$ \\
$\int\rho\,\mathrm{d}\rr$ from \file{CHGCAR} & $297$ & $297.0000$ \\
Spectral resampling, band-limited & exact & $1.3\times10^{-15}$ \\
Resampling, mean preserved & exact & $5.6\times10^{-17}$ \\
Madelung constant (rock salt) & $1.7475645946$ & $1.74756459$ \\
\file{EXTCAR} vs VASP, tabulated & --- & $2.1\times10^{-5}$ rel.\ $L^2$ \\
\bottomrule
\end{tabular}
\end{center}

The electron count is the sharpest of these: it exercises the \file{CHGCAR}
parser, the $\Omega$ convention and the grid metric simultaneously, and returns
the exact integer.

\section{Model 1: \texorpdfstring{$\vext \mapsto \rho$}{Vext to rho}}

Leave-one-out over the three materials; metrics on the held-out material in
physical units ($e/\text{\AA}^3$).

\begin{center}
\begin{tabular}{@{}lrrrr@{}}
\toprule
Held out & MAE & RMSE & rel.\ $L^2$ & $R^2$ \\
\midrule
struct\_000 & $0.01667$ & $0.03097$ & $0.0318$ & $0.9983$ \\
struct\_001 & $0.02433$ & $0.03984$ & $0.0405$ & $0.9973$ \\
struct\_002 & $0.02423$ & $0.04112$ & $0.0418$ & $0.9971$ \\
\midrule
\textbf{mean} & $\mathbf{0.02174}$ & $\mathbf{0.03731}$ & $\mathbf{0.0380}$ & $\mathbf{0.9976}$ \\
\midrule
baseline: mean density & $0.5199$ & $0.7589$ & $0.7779$ & $0.0000$ \\
\bottomrule
\end{tabular}
\end{center}

The learned map is a factor of $20$ better than predicting a constant density.
The electron count, which nothing in the objective constrains, is recovered to
$1.1$--$2.2\,\%$ --- precisely the single global degree of freedom that the
charge-normalisation head of Section~\ref{sec:pauli-head} would fix exactly and
for free.

\section{Model 2: \texorpdfstring{$\rho \mapsto \tau$}{rho to tau}}

Same protocol, in $\mathrm{eV}/\text{\AA}^3$, with the Pauli head of
Eq.~\eqref{eq:paulihead} enabled. The analytic orbital-free functionals of
Chapter~\ref{ch:ofdft} are evaluated on the same input as baselines.

\begin{center}
\begin{tabular}{@{}lrrrr@{}}
\toprule
Predictor & MAE & RMSE & rel.\ $L^2$ & $R^2$ \\
\midrule
\textbf{FNO (learned)} & $\mathbf{0.6600}$ & $\mathbf{1.2002}$ & $\mathbf{0.0620}$ & $\mathbf{0.9924}$ \\
\midrule
von Weizsäcker & $8.99$ & $14.22$ & $0.7367$ & $0.0150$ \\
Thomas--Fermi & $9.38$ & $25.99$ & $1.3454$ & $-2.2908$ \\
TF $+$ vW$/9$ & $9.35$ & $26.18$ & $1.3554$ & $-2.3395$ \\
\bottomrule
\end{tabular}
\end{center}

The learned functional is better than the best classical approximation by a
factor of about $12$ in relative $L^2$.

\begin{pnote}
Thomas--Fermi attains a \emph{negative} coefficient of determination: as a
pointwise predictor of $\tau$ it is worse than the constant mean. This is not a
defect of the implementation but the expected failure of a uniform-electron-gas
limit applied to a $d$-band metal, whose core and $d$-shell regions are about as
far from uniform as matter gets. It is the gap that motivates a learned
functional.
\end{pnote}

\subsection{Integrated kinetic energy}

The quantity that enters a total energy is $\Ts=\int\tau\,\mathrm{d}\rr$, not
$\tau$ pointwise:

\begin{center}
\begin{tabular}{@{}lrrr@{}}
\toprule
Held out & predicted [eV] & reference [eV] & error \\
\midrule
struct\_000 & $6277.16$ & $6207.07$ & $1.13\,\%$ \\
struct\_001 & $6198.07$ & $6214.80$ & $0.27\,\%$ \\
struct\_002 & $6180.76$ & $6217.06$ & $0.58\,\%$ \\
\bottomrule
\end{tabular}
\end{center}

\subsection{Constraint satisfaction}

With the head enabled, $\tau \ge \tauvW[\rho]$ held at every one of the
$\sim3\times10^{4}$ grid points on all three folds, with minimum margins of
$1.045$, $0.950$ and $1.024\ \mathrm{eV}/\text{\AA}^3$. The margins are the
informative number: the model is not sitting on the bound but predicting a Pauli
term comfortably above it.

The fitted scale $s$ came out at $5.22$, $5.02$ and $5.01\
\mathrm{eV}/\text{\AA}^3$ across folds --- consistent with the median $\tauP$
measured directly from the data ($4.81$--$5.21$), which is a check that
\opt{fit\_pauli\_scale} estimates what it claims to.

Section~\ref{sec:headresult} gives the matched comparison against an
unconstrained backbone: the head improves every fold and every metric, by
$18.3\,\%$ in mean relative $L^2$, while training $17.8\,\%$ \emph{faster}.

\section{The inference pipeline}

Chaining both models turns a bare crystal geometry into predicted $\rho$ and
$\tau$ with no wavefunctions and no self-consistency cycle:
\begin{equation}
  \{\text{\file{POSCAR}}, \text{\file{INCAR}}, \text{\file{POTCAR}}\}
  \;\xrightarrow{\ \text{analytic}\ }\; \vext
  \;\xrightarrow{\ \text{FNO 1}\ }\; \rho
  \;\xrightarrow{\ \text{FNO 2}\ }\; \tau .
\end{equation}
Only the two field-to-field maps are learned; the first step is closed-form.

The importance of Chapter~\ref{ch:vasp}'s correction shows up here most
sharply. Evaluating on struct\_002 with models that never saw it:

\begin{center}
\begin{tabular}{@{}lrr@{}}
\toprule
Field & Gaussian $\vext$ & \textbf{tabulated $\vext$} \\
\midrule
\file{EXTCAR} & $0.6645$ & $\mathbf{0.0014}$ \\
\file{CHGCAR} & $0.4482$ & $\mathbf{0.0647}$ \\
\file{TAUCAR} & $0.2456$ & $\mathbf{0.1190}$ \\
\bottomrule
\end{tabular}
\end{center}

(relative $L^2$ against the reference files.) The models were always
\emph{trained} on VASP's own \file{EXTCAR}; the error was entirely in the
inference path, which had been feeding them a potential from a different
distribution. Correcting the analytic step alone improved the predicted density
sevenfold, with no retraining.

\begin{ptip}
All predicted fields are written in \file{CHGCAR} format.
\opt{--compare} additionally brings reference files onto the prediction
grid by spectral truncation --- the exact restriction --- rather than failing on
a shape mismatch.
\end{ptip}

\section{Hardware}

The full run --- two tasks, three folds each, $200$ epochs, spectral
downsampling, figures and checkpoints --- completes in roughly $20$ minutes on
Apple Silicon via MPS. Forward-pass throughput is given in
Section~\ref{sec:hardware}; the advantage grows with grid size, from
$2.2\times$ at $32^3$ to $5.7\times$ at $64^3$, because kernel-launch overhead
dominates at small sizes while the FFTs dominate at large ones.

\section{What these numbers do and do not support}

\begin{pwarn}
With three closely related structures of a single element, leave-one-out
measures \textbf{interpolation between nearby geometries of platinum}. It
demonstrates that the pipeline is correct, that the architecture has sufficient
capacity, and that the learned kinetic functional beats the classical
approximations on this system. It says \textbf{nothing} about transfer to other
chemistry, other coordination environments, or other cutoffs. The training log
carries this caveat in its own footer so that it travels with the numbers.
\end{pwarn}

What \emph{is} established independently of data-set size:

\begin{itemize}
  \item the field machinery is verified against closed-form physics
    (Poisson, Madelung, electron count) to machine precision;
  \item the \file{EXTCAR} reconstruction matches a reference implementation to
    $2\times10^{-5}$;
  \item one set of weights serves three different grid shapes, with
    CPU/accelerator agreement at $10^{-6}$;
  \item the von Weizsäcker bound is structural, verified at random
    initialisation and under adversarial optimisation;
  \item the classical functionals' failure on this system is quantified rather
    than asserted.
\end{itemize}

The obvious next step is more materials. The ingestion layer, the shape-bucketed
sampler and the configuration system were all built for a heterogeneous data
set; the binding constraint is now the data, not the code.
